\section*{September 7, 2026: Require complete component measurements}
\textit{Agent-drafted entry; Jacob approved the assessment-snapshot selection rule.}

Selecting the preferred assessment separately for each building card can produce
a total that never appears in any assessment year. Four mechanically retained
multiple-card candidates had this defect. Another candidate selected an old
pre-redevelopment report for one card and a new-building report for another,
omitting part of the new construction. These are problems with assembling source
records, before any density regression is run.

Jacob approved requiring a complete snapshot under the existing 2022/2025/later
priority. The card producer now selects one assessment for each proposed
construction-year episode. The snapshot must contain exactly the selected
component cards, with positive units and building areas and one positive parcel
land value. An extra card is a membership conflict, not permission to choose a
convenient subset. The project producer requires support from one assessment year
before adding the components. The later evidence reader consumes that selection
instead of constructing a second one.

The candidate inventory still contains 27,831 records. Construction years, parcel
membership, land values, and preliminary boundary distances are unchanged. Three
previously mechanically retained candidates now require reconciliation: the
Washington Boulevard records never provide the selected components together;
Ingleside has an extra rebuilt component; and another group combines retired and
reclassified cards. The intermediate review count rises from 190 to 193.
No new sample exclusion or case-specific ledger entry was introduced.

Twelve unit totals and twelve building-area totals become missing because their
components lack complete support. Three other floor-area totals now come from
complete snapshots: 3,934 becomes 3,976 square feet; 3,302 becomes 3,289; and a
partial six-card sum of 4,284 becomes 12,744. The six-card project already required
review because its cards report different construction years. These are changes
to Assessor candidate measurements, not adopted final density outcomes.

Snapshot consistency does not establish physical identity. The existing review
for 42 E 90th Street identifies one house despite two cards appearing together
in 2009. Its physical-identity correction must still enter downstream. Card
renumbering, redevelopment, and subdivision therefore remain the next membership
questions. The frozen paper input was not replaced and no new estimates were run.

\small
\textit{Reproduction and scope.} The actual construction Make build regenerates
\path{residential_multicard_cards.csv} and
\path{residential_assessor_project_candidates.csv} and their standard reports.
The comparison is against revision \texttt{36094f8} using the same pinned sources.
All 2,209 card rows with a selected complete snapshot match their source records;
no selected construction-year group mixes assessment years. A downstream Make
replay holds later project, permit, and geographic inputs fixed, so it verifies
the affected readers rather than a complete final-sample rebuild. Comparison
hashes are recorded in the release audit's \path{restored_stage_checks.csv}.
\normalsize
